\section{A velocity circle with controlled turns}
\label{sec:velocity-loop}

We now construct the leading oscillation for
Proposition~\ref{prop:primitive-addition}.  Throughout this section its data
$F,D,a,x,y,n$ and the later curves are fixed.  The basic choice is to
oscillate in the plane perpendicular to both $G_y$ and $G_{yy}$, for a
slow map $G$ near $F$.  Orthogonality to the second derivative will remove
the first error in the $yy$ metric coefficient.  A separate choice at the
two turning points of the velocity path will preserve the later boundary
conditions.  For a $2\pi$-periodic function $b$ we write
$\avg b=(2\pi)^{-1}\int_0^{2\pi}b(\theta)\,d\theta$.
For a nonzero vector $b$, write $\widehat b=b/|b|$.

The use of a periodic velocity to add a rank-one metric tensor belongs
to the oscillatory constructions of Nash and Kuiper
\cite{Nash1954,Kuiper1955}; see also \cite[\S4]{Choquet1958} and
\cite[\S\S3.2--3.3]{Gromov2015}.
The second-derivative constraint on the velocity plane and the behavior
prescribed at its turns are part of the construction proved here.

\subsection{The plane and the leading metric}

For a 2-jet near that of $F$, define
\begin{equation}
 Y=G_y,\qquad C=G_{yy},\qquad
 P=\operatorname{span}(Y,C)^\perp,\qquad
 v=PG_x,\qquad R_0=(|v|^2+a^2)^{1/2}.
 \label{eq:velocity-plane}
\end{equation}
We use $P$ also for orthogonal projection to this plane.  The following
facts justify these definitions and the frame we will use.

\begin{lemma}
\label{lem:velocity-plane}
On a neighborhood of $\overline D$ and a sufficiently small neighborhood
of the 2-jet section of $F$, the vectors $Y,C$ are independent, $P$ is a
smooth 2-plane, and $R_0$ is bounded below by a positive constant.
Projection to $P$ restricts to an isomorphism on the normal plane of $G$.
There is a smooth orthonormal frame $(e_1,e_2)$ of $P$, depending smoothly
on these jets, such that $e_1=v/|v|$ on a collar of $C_0$.
\end{lemma}

\begin{proof}
In the coordinates $(x,y)$,
\[
 \Hess_{\gamma(G)}x(\partial_y,\partial_y)=-\Gamma^x_{yy},\qquad
 C=\Gamma^x_{yy}G_x+\Gamma^y_{yy}G_y+B_G(\partial_y,\partial_y).
\]
By \eqref{eq:primitive-hessian}, $\Gamma^x_{yy}<0$ at $F$ on
$\overline D$, and this inequality persists on the indicated
neighborhoods.  The tangential projections of $Y,C$ are therefore
independent.  If a normal vector belongs to $\operatorname{span}(Y,C)$,
its zero tangential projection forces both coefficients to vanish.
Projection from the normal plane to $P$ is consequently an isomorphism.
Projecting the displayed decomposition of $C$ gives
\begin{equation}
 P B_G(\partial_y,\partial_y)=-\Gamma^x_{yy}v.
 \label{eq:projected-yy-form}
\end{equation}
At $F$, $A_0\ne0$ on $C_0$, so $v\ne0$ on a collar.  Moreover $v$ and
$Pn$ are not opposite rays there, by
\eqref{eq:boundary-invariants}, \eqref{eq:projected-yy-form}, and the
injectivity of projection on the old normal plane.

The vector $Pn$ is nonzero over the whole disk at $F$, and remains so for
nearby jets.  On a smaller collar set $e_1=\widehat v$, and in the
interior set $e_1=\widehat{Pn}$.  In the intervening collar, normalize a
smooth convex interpolation of these two non-antipodal vectors.
Orient $P$ using the ordered pair $(Y,C)$ and the orientation of $\R^4$,
and complete $e_1$ to $e_2$ with that orientation.  These choices are
local to this coordinate disk.

Finally $R_0>0$ because $a>0$ in $D$ and $v\ne0$ near its boundary.
Compactness gives a positive lower bound, which persists after shrinking
the jet neighborhood and the exterior neighborhood.
\end{proof}

For an angular function $\alpha$ to be chosen, set
\begin{equation}
 p=\cos\alpha\,e_1+\sin\alpha\,e_2,\qquad
 V=R_0p,\qquad w=V-v,\qquad X=G_x+w.
 \label{eq:leading-velocity}
\end{equation}
Since $G_x-v\in\operatorname{span}(Y,C)$ and $V,v\in P$,
\begin{equation}
 |X|^2=|G_x|^2+a^2,\qquad X\cdot Y=G_x\cdot G_y.
 \label{eq:leading-metric}
\end{equation}
Thus the pair $(X,Y)$ has exactly the required leading metric.  What
remains is to make the periodic mean of $w$ zero, while controlling the
places where the velocity stops along its circle.

\subsection{A prescribed mean and smooth vanishing}

Initially use an angular parameter $t\in\R/(2\pi\mathbb Z)$ and angles
\begin{equation}
 \alpha=h_{\mathrm{turn}}+A\cos t.
 \label{eq:unreparametrized-angle}
\end{equation}
Here $h_{\mathrm{turn}}$ will be a smooth function supported in a fixed
compact subset of $D$ and small in value.  We first choose $A$, independently
of that function.

On a sufficiently small collar, $e_1=\widehat v$ and $a/|v|$ is small.
There is a nonnegative smooth amplitude $A_{\mathrm{sm}}$ satisfying
\begin{equation}
 \frac1{2\pi}\int_0^{2\pi}\cos(A_{\mathrm{sm}}\cos t)\,dt
 =\frac{|v|}{R_0}.
 \label{eq:small-arc-mean}
\end{equation}
Indeed the left side is $1-A_{\mathrm{sm}}^2/4+O(A_{\mathrm{sm}}^4)$.
The implicit function theorem applied to the squared amplitude gives
\begin{equation}
 A_{\mathrm{sm}}=\frac a{|v|}\,
 c\left(\frac{a^2}{|v|^2}\right),\qquad c(0)=\sqrt2,
 \label{eq:smooth-small-amplitude}
\end{equation}
where $c$ is positive and smooth.  Thus $A_{\mathrm{sm}}$ is smooth and
flat on $C_0$, and it is zero on the exterior.  When
$h_{\mathrm{turn}}=0$, the sine average vanishes as well, so the
unweighted mean of $R_0p$ is exactly $v$ on this collar.

Farther inward, while still in a region where $e_1=\widehat v$ and the
small formula is defined, interpolate $A$ from $A_{\mathrm{sm}}$ to a
fixed number larger than $\pi$, keeping $A\ge A_{\mathrm{sm}}$.
Use that constant on the remaining interior.  For $a>0$ the point
$v/R_0$ lies strictly inside the convex hull of the path $p$ when
$h_{\mathrm{turn}}=0$.  On the small arc this follows because its
unweighted mean lies strictly between the endpoint chord and the point
$(1,0)$ in the $(e_1,e_2)$ coordinates.  Increasing $A$ retains this arc;
when $A>\pi$ the path covers the circle and $|v|/R_0<1$.
On any fixed compact subset of the interior, this strict membership
persists for sufficiently small $\|h_{\mathrm{turn}}\|_{C^0}$, uniformly
on a compact neighborhood of the permitted jet data.

We use the following elementary fact to adjust the mean.

\begin{lemma}[Smooth densities for interior averages]
\label{lem:loop-density}
Let a smooth family of circle-valued paths $p_b(t)\in\R^2$ and targets
$c_b\in\R^2$ be parametrized by a finite-dimensional parameter $b$.
On a compact parameter set, suppose each $c_b$ lies in the interior of
the convex hull of $p_b$.  Then, on a neighborhood of that set, there
are smoothly varying strictly positive densities $\rho_b(t)$ such that
\[
 \avg_t\rho_b=1,\qquad \avg_t(\rho_b p_b)=c_b.
\]
These densities may be patched with any already prescribed positive
solutions on an overlapping parameter neighborhood.
\end{lemma}

\begin{proof}
At a fixed parameter choose finitely many path values whose convex hull
contains the target in its interior.  Choose strictly positive weights
with this barycenter and total mass one.  Replacing the point masses by
sufficiently narrow smooth bumps, and adding a sufficiently small
positive constant density, preserves the interior condition.
The three linear constraints, for mass and the two coordinates, have
rank three.  A fixed right inverse of their coefficient matrix gives
nearby positive weights solving the constraints smoothly as the
parameter varies.  A finite parameter cover and a smooth partition of
unity patch these local solutions: the constraints are linear in the
density and positivity is preserved by convex combinations.  The same
argument patches with an already prescribed solution.
\end{proof}

Apply the lemma in the frame $(e_1,e_2)$, with $c=v/R_0$, and use
$\rho=1$ on a smaller collar where
\eqref{eq:small-arc-mean} applies and $h_{\mathrm{turn}}=0$.
Here the parameters are position and the $2$-jet of $G$.  The frame
identifies all the velocity planes with $\R^2$, and the remaining
parameter set is compact and lies in $a>0$.  The interior hypothesis
therefore has the required uniform margin, and the lemma patches a
smooth density with the prescribed collar solution.  Put
\[
 \theta=\int_0^t\rho(s)\,ds
\]
and replace $t$ by the smooth inverse of this orientation-preserving
circle diffeomorphism.  We henceforth use $\theta$ as periodic parameter.
Its uniform mean satisfies
\begin{equation}
 \avg_\theta V=v,\qquad \avg_\theta w=0.
 \label{eq:velocity-mean}
\end{equation}
All data are smooth functions of position and the 2-jet of $G$.
On the exterior $A=0$, $h_{\mathrm{turn}}=0$, and $\rho=1$, so $V=v$.
Equations \eqref{eq:smooth-small-amplitude} and
\eqref{eq:small-arc-mean} show that $w$ is smooth across $C_0$ and
vanishes there with every spatial derivative.

In the interior $A>0$.  There are exactly two zeros of $\alpha_\theta$,
corresponding to $t=0,\pi$.  At these turns,
\begin{equation}
 \alpha=h_{\mathrm{turn}}\pm A,
 \qquad \alpha_y=(h_{\mathrm{turn}})_y\pm A_y.
 \label{eq:turn-derivative}
\end{equation}
In this formula spatial differentiation holds the new parameter
$\theta$ fixed.  The additional chain-rule term from the inverse
reparametrization is zero because it is multiplied by $\sin t$.
Figure~\ref{fig:velocity-collar} separates the velocity circle from its
angular parametrization.  The two stationary angles persist under the
positive reparametrization; their transverse derivatives are the quantities
we will prescribe along the later curves.

\begin{figure}[htbp]
\centering
\begin{tikzpicture}[scale=.88,>=Latex,font=\small]
 \begin{scope}
  \draw[->,gray] (-2.35,0)--(2.65,0);
  \draw[->,gray] (0,-2.25)--(0,2.3);
  \draw[thick,blue!60!black] (0,0) circle (1.85);
  \fill (0,0) circle (1.6pt) node[below left] {$0$};
  \fill (.75,0) circle (2pt) node[below right=3pt] {$v$};
  \coordinate (V) at (60:1.85);
  \draw[->,thick] (0,0)--(V) node[midway,left] {$R_0$};
  \draw[->,thick,blue!70!black] (.75,0)--(V)
    node[midway,right=3pt] {$V-v$};
  \fill (V) circle (2pt) node[above right] {$V$};
  \node at (0,-2.85) {Velocity circle};
 \end{scope}
 \begin{scope}[xshift=4cm]
  \draw[->,gray] (0,-1.65)--(0,1.95) node[above] {$\alpha$};
  \draw[->,gray] (0,0)--(3.65,0) node[right] {$t$};
  \draw[dashed,gray!60] (0,1.2)--(3.3,1.2);
  \draw[dashed,gray!60] (0,-1.2)--(3.3,-1.2);
  \node[left] at (0,1.2) {$h_{\mathrm{turn}}+A$};
  \node[left] at (0,0) {$h_{\mathrm{turn}}$};
  \node[left] at (0,-1.2) {$h_{\mathrm{turn}}-A$};
  \draw[thick,blue!60!black,domain=0:360,samples=81]
    plot ({3.3*\x/360},{1.2*cos(\x)});
  \fill[blue!60!black] (0,1.2) circle (2pt);
  \fill[blue!60!black] (1.65,-1.2) circle (2pt);
  \fill[blue!60!black] (3.3,1.2) circle (2pt);
  \node[above right] at (0,1.25) {turn};
  \node[below] at (1.65,-1.3) {turn};
  \node[below right] at (0,0) {$0$};
  \draw (1.65,.06)--(1.65,-.06) node[below] {$\pi$};
  \draw (3.3,.06)--(3.3,-.06) node[below] {$2\pi$};
  \node at (1.65,-2.85) {Angle before reparametrization};
 \end{scope}
 \begin{scope}[xshift=10.8cm]
  \fill[gray!12] (0,0) circle (1.95);
  \fill[white] (0,0) circle (1.56);
  \draw[thick] (0,0) circle (1.95);
  \draw[dashed,gray] (0,0) circle (1.56);
  \draw[thick,blue!60!black] (-2.3,-1.15) .. controls (-.7,-.7) and (.4,.8) .. (2.2,.95);
  \draw[thick,red!65!black] (-1.8,1.85) .. controls (-.9,.3) and (.9,-.2) .. (2.25,-1.25);
  \node[blue!60!black,above right] at (2.05,.95) {$C_i$};
  \node[red!65!black,below right] at (2.1,-1.18) {$C_l$};
  \node[above] at (1.12,1.63) {$C_0$};
  \node at (-1.4,-2.15) {collar};
  \draw[->,thin] (-1.05,-2.04)--(-.65,-1.72);
  \node at (.9,.8) {$D$};
  \node at (0,-2.85) {Support collar};
 \end{scope}
\end{tikzpicture}
\caption{Velocity circle, its angular parametrization, and the support collar.
The circle in $P=\operatorname{span}(Y,C)^\perp$ has radius
$R_0=(|v|^2+a^2)^{1/2}$; $V-v$ is the correction vector.
The middle panel plots $\alpha=h_{\mathrm{turn}}+A\cos t$ before the
positive reparametrization that gives mean velocity $v$.
It has two turns per period, at $t=0\equiv2\pi$ and $t=\pi$.
There $\alpha_y=(h_{\mathrm{turn}})_y\pm A_y$, even after
reparametrization, so a large transverse derivative can control both turns.
The angle graph also applies when $A>\pi$ and the velocity path wraps
around the circle.  As $a$ vanishes at $C_0$, the path contracts smoothly
to $v$, while the circle retains radius $|v|>0$.
The exact immersions approach the old map in $C^2$ outside $D$
as the fast scale tends to zero.}
\label{fig:velocity-collar}
\end{figure}


\subsection{A small function with a prescribed transverse derivative}

We will need a large value of $(h_{\mathrm{turn}})_y$ along the later
curves away from the boundary, while keeping $h_{\mathrm{turn}}$ small
in value.  The finite-tangency hypothesis of
Proposition~\ref{prop:primitive-addition} gives exactly this freedom.

\begin{lemma}
\label{lem:small-value-derivative}
Let $K$ be a compact subset of a finite union of smooth embedded curves
in the interior of a coordinate disk.  Suppose the curves have only
finitely many intersections and only finitely many tangencies to
$\partial_y$ near $K$, with no shared curve arcs.  Given $B>0$ and
$\epsilon>0$, there is $h\in C_c^\infty(D)$ such that
\[
 \|h\|_{C^0}<\epsilon,\qquad h_y=B\quad\text{on }K.
\]
The support may be kept in any fixed precompact open neighborhood of $K$.
\end{lemma}

\begin{proof}
Place the finitely many intersections and vertical tangencies in
pairwise disjoint small disks centered at points $p$.  Choose bumps
$\chi_p$ equal to one on smaller disks and supported in the larger ones.
The function
\[
 h^{\mathrm{pre}}=\sum_p\chi_p B(y-y(p))
\]
has arbitrarily small sup norm when the disks are sufficiently small,
and $(h^{\mathrm{pre}})_y=B$ on the smaller disks.
Choose still smaller cores.  Outside those cores, the relevant compact
curve set is covered by finitely many patches containing just one
curve.  In such a patch choose a defining function $f_j$ with
$(f_j)_y\ne0$; the latter condition expresses transversality to
$\partial_y$.  Choose smooth patch cutoffs $\chi_j$ whose sum is one
on the part of $K$ where $(h^{\mathrm{pre}})_y\ne B$.
Add the functions
\[
 h_j=\chi_j f_j
       \frac{B-(h^{\mathrm{pre}})_y}{(f_j)_y}
       \eta_j(f_j/\delta_j),
\]
where $\eta_j$ is one near zero and has compact support.
On the curve $f_j=0$, differentiation gives exactly
\[
 (h_j)_y=\chi_j\bigl(B-(h^{\mathrm{pre}})_y\bigr).
\]
All derivatives of the cutoffs in this identity have a factor $f_j$
and therefore vanish there.  Thus
$h=h^{\mathrm{pre}}+\sum_jh_j$ has $h_y=B$ on $K$, including the
transition annuli of the special-point bumps.  Each quotient is bounded
on its fixed patch, so taking the finitely many $\delta_j$ sufficiently
small makes the additional sup norms as small as desired.  The supports
stay within the prescribed neighborhood.
\end{proof}

\subsection{Uniform normal bounds and the choice at turns}

For the rest of this section evaluate the slow data at $G=F$.
Let $S_0$ be the length of the projection of $C$ to
$\operatorname{span}(X,Y)^\perp$, and let $n_0^*$ be its unit direction.
This projection never vanishes: if $C$ belonged to
$\operatorname{span}(X,Y)$, applying $P$ would force its $X$ coefficient
to be zero since $PX=V\ne0$, contradicting independence of $C,Y$.
Compactness gives positive lower and upper bounds for $S_0$, uniformly
over all angular values.  A complementary unit normal is
\[
 m_0^*=-\sin\alpha\,e_1+\cos\alpha\,e_2.
\]
It is perpendicular to $X,Y,C$, and hence also to $n_0^*$.  Set
\begin{equation}
 D_0^*=\langle X_y,n_0^*\rangle,
 \qquad N_0^*=\langle X_y,m_0^*\rangle.
 \label{eq:limiting-mixed-components}
\end{equation}
The derivative here holds the periodic parameter fixed.  Directly,
\[
 X_y=(F_x-v)_y+(R_0)_yp+
 R_0\bigl(\cos\alpha(e_1)_y+\sin\alpha(e_2)_y
                    +\alpha_y m_0^*\bigr).
\]
The term containing $\alpha_y$ drops out of $D_0^*$, so
\begin{equation}
 |D_0^*|\le C_0,
 \qquad N_0^*=R_0\alpha_y+E_0,\quad |E_0|\le C_0,
 \label{eq:derivative-independent-bounds}
\end{equation}
where $C_0$ is independent of $h_{\mathrm{turn}}$, its derivatives,
and the reparametrization.  These estimates hold over all angular
values.  The same independence holds for the bounds on $S_0$ and
$1/S_0$.

Write
\[
 k_+=K_{\gamma_+}\det(\gamma_+)
\]
with determinant in $(x,y)$ coordinates, and write
$Y_i=b_i\partial_x+c_i\partial_y$ on the disk.  At a crossing
$C_i\cap C_0$, the loop is constant and $X=F_x$.  Since $a$ is flat,
the prescribed metric $\gamma_+$ has the old curvature there.
Consequently the strict condition $Q_{0i}(F)>0$ is exactly
\begin{equation}
 (D_0^*b_i+S_0c_i)^2+k_+b_i^2>0.
 \label{eq:collar-inequality}
\end{equation}
It remains valid on the later curves in a sufficiently thin collar for
all path angles.  Indeed $A_{\mathrm{sm}}$ tends to zero at the boundary,
and $D_0^*$ contains no angular derivative.  Compactness confines the
later curve portions of a thin collar to neighborhoods of their
finitely many crossings with $C_0$.  The collar can also avoid every
later mutual crossing, because there are no triple intersections.
The same inequality holds if the path angles are translated by any
sufficiently small value of $h_{\mathrm{turn}}$.  Thus this collar
estimate remains available even if the eventual interior support
neighborhood of $h_{\mathrm{turn}}$ overlaps its inner part.

Let $K$ be the compact union of the later curve portions in the disk
outside this collar.  By Lemma~\ref{lem:small-value-derivative},
we may prescribe $(h_{\mathrm{turn}})_y=B_*$ on $K$ with arbitrarily
small values of $h_{\mathrm{turn}}$.  The support can lie in a fixed compact neighborhood of $K$ inside $D$,
while a smaller boundary collar remains unchanged.  On any overlap
with the collar estimate, smallness in value preserves that estimate.
For any fixed threshold $T$, choose
\begin{equation}
 B_* > \|A_y\|_{C^0}
        +\frac{T+C_0+1}{\min R_0}.
 \label{eq:turn-size-choice}
\end{equation}
Equations \eqref{eq:turn-derivative} and
\eqref{eq:derivative-independent-bounds} then give
$N_0^*>T+1$ at both turns on $K$.  After making the values of
$h_{\mathrm{turn}}$ small enough to retain convex-hull membership,
choose the positive density as above.  That reparametrization does not
change the turn derivative identity.

We summarize the construction in the form needed by the next section.

\begin{lemma}[Velocity loop]
\label{lem:velocity-loop}
Under the hypotheses of Proposition~\ref{prop:primitive-addition}, fix
any finite threshold $T$.  There is a neighborhood of the 2-jet section
of $F$ over $\overline D$ and a small exterior neighborhood, on which
one can choose smooth $2\pi$-periodic vectors $V,w$ as in
\eqref{eq:velocity-plane}--\eqref{eq:leading-velocity}, with
\[
 |V|=R_0>0,\qquad \avg V=v,\qquad \avg w=0.
\]
They depend smoothly on position and the 2-jet; $w$ extends smoothly by
zero outside $D$.  The leading derivatives $(X,Y)$ have metric
$\gamma(G)+a^2dx^2$.  At $G=F$, the limiting normal data defined in
\eqref{eq:limiting-mixed-components} have uniform bounds
\eqref{eq:derivative-independent-bounds}, a positive lower bound for
$S_0$, and the collar inequality \eqref{eq:collar-inequality}.
On the remaining compact portions of the later curves, the only zeros
of $\alpha_\theta$ are the two turns, and there $|N_0^*|>T$.
The bounds for $S_0,1/S_0,D_0^*$ and the error $E_0$ can be fixed before
choosing $T$ and $B_*$.
\end{lemma}

All choices in this section are now fixed before the fast oscillation
scale is sent to zero.  In particular their higher derivatives may be
large but are finite.  The next section will use this fixed smooth
loop to produce exact immersions.  We will choose $T$ from explicit
crossing bounds in Section~\ref{sec:geometry-preservation}; those bounds
use only the quantities just shown to be independent of $T$.
